\ifcsname LRAComputationalLinearAlgebraCapstoneRendered\endcsname
\else
\expandafter\gdef\csname LRAComputationalLinearAlgebraCapstoneRendered\endcsname{}

\newpage
\phantomsection
\label{cap:computational-linear-algebra}

\begin{tcolorbox}[
  colback=gray!6,
  colframe=gray!40,
  arc=2pt,
  left=8pt, right=8pt, top=6pt, bottom=6pt,
  title={\small\textbf{Capstone Problem}},
  fonttitle=\small\bfseries
]
\textbf{Problem.}
TODO: state one synthesizing problem for Computational Linear Algebra, using only concepts at or
before this chapter in the registry.
\end{tcolorbox}

\begin{remark*}[Dependency ceiling]
The capstone may use only results routed at or before this chapter.
\end{remark*}

% =========================================================
% Migrated exercises from exercises-1.tex
% =========================================================

% =========================================================
% Computational Linear Algebra --- Exercises
% Focus: NURBS/B\'{e}zier readiness --- computational, no proofs
% =========================================================

\subsection*{Exercises: Computational Linear Algebra}

% ---------------------------------------------------------
\subsubsection*{Vectors and Linear Combinations}

\begin{enumerate}

\item Compute the following:
\begin{enumerate}
  \item $(2, -1, 3) + (-1, 4, -2)$
  \item $3(1, 0, -2) - 2(-1, 1, 1)$
  \item $\|(3, -4)\|$, $\|(1, 2, 2)\|$, and $\|(0, 0, 0, 5)\|$
\end{enumerate}

\item Write $(5, -3, 7)$ as a linear combination of the standard basis
vectors $\mathbf{e}_1, \mathbf{e}_2, \mathbf{e}_3 \in \mathbb{R}^3$.

\item Determine whether $(1, 2, 3)$ is a linear combination of
$(1, 1, 0)$ and $(0, 1, 1)$.
If yes, find the coefficients. If no, explain why not.

\item The control points of a quadratic B\'{e}zier curve are
$P_0 = (0,0)$, $P_1 = (1,2)$, $P_2 = (3,0)$.
The curve is $C(u) = (1-u)^2 P_0 + 2u(1-u)P_1 + u^2 P_2$.
\begin{enumerate}
  \item Compute $C(0)$, $C(1/2)$, and $C(1)$.
  \item Verify that $C(1/2)$ is a convex combination of $P_0, P_1, P_2$
        (check that the weights are non-negative and sum to 1).
  \item Verify that $C(1/2)$ lies inside the triangle formed by the control points.
\end{enumerate}

\item Let $P_0 = (0,0)$, $P_1 = (2,3)$, $P_2 = (4,0)$ be control points.
The de Casteljau algorithm at $u = 1/3$ computes:
\[
  P_0^1 = (1-u)P_0 + uP_1, \quad P_1^1 = (1-u)P_1 + uP_2, \quad
  C(u) = (1-u)P_0^1 + uP_1^1.
\]
Carry out this computation and find $C(1/3)$.

% ---------------------------------------------------------
\subsubsection*{Dot Products and Lengths}

\item Compute the dot product and the angle between the following pairs of vectors.
State whether they are orthogonal, acute, or obtuse.
\begin{enumerate}
  \item $\mathbf{v} = (1,0,0)$ and $\mathbf{w} = (0,1,0)$
  \item $\mathbf{v} = (1,2,3)$ and $\mathbf{w} = (4,-2,0)$
  \item $\mathbf{v} = (1,1)$ and $\mathbf{w} = (1,-1)$
\end{enumerate}

\item Find the unit vector in the direction of $(3, -4, 0)$.
Then find the unit vector in the direction of $(1, 1, 1)$.

\item The tangent vector to the quadratic B\'{e}zier curve at $u=0$ is
$C'(0) = 2(P_1 - P_0)$, and at $u=1$ is $C'(1) = 2(P_2 - P_1)$.
Using the control points $P_0 = (0,0)$, $P_1 = (2,1)$, $P_2 = (4,0)$:
\begin{enumerate}
  \item Compute $C'(0)$ and $C'(1)$.
  \item Find the angle between the two endpoint tangent vectors.
  \item Find the unit tangent vectors at each endpoint.
\end{enumerate}

% ---------------------------------------------------------
\subsubsection*{Matrices and Matrix Operations}

\item Compute the following matrix products, or state why they are not defined.
\begin{enumerate}
  \item $\begin{pmatrix} 1 & 2 \\ 3 & 4 \end{pmatrix}
         \begin{pmatrix} 0 & 1 \\ 1 & 0 \end{pmatrix}$

  \item $\begin{pmatrix} 1 & 0 & 2 \\ 0 & 1 & -1 \end{pmatrix}
         \begin{pmatrix} 2 \\ -1 \\ 3 \end{pmatrix}$

  \item $\begin{pmatrix} 2 \\ -1 \\ 3 \end{pmatrix}
         \begin{pmatrix} 1 & 0 & -1 \end{pmatrix}$
\end{enumerate}

\item The quadratic B\'{e}zier curve $C(u) = \sum_{i=0}^2 B_{i,2}(u) P_i$ can be
written in matrix form as
\[
  C(u) = \begin{pmatrix} 1 & u & u^2 \end{pmatrix}
         \begin{pmatrix} 1 & 0 & 0 \\ -2 & 2 & 0 \\ 1 & -2 & 1 \end{pmatrix}
         \begin{pmatrix} P_0 \\ P_1 \\ P_2 \end{pmatrix}.
\]
Using $P_0 = (0,0)$, $P_1 = (1,2)$, $P_2 = (2,0)$, compute $C(1/2)$
by carrying out the matrix products in order from right to left.

\item Let
$A = \begin{pmatrix} 2 & 1 \\ 0 & 3 \end{pmatrix}$
and
$B = \begin{pmatrix} 1 & -1 \\ 2 & 0 \end{pmatrix}$.
Compute $AB$, $BA$, and confirm that $AB \neq BA$.
Also compute $(AB)^T$ and $B^T A^T$, and confirm they are equal.

% ---------------------------------------------------------
\subsubsection*{Inverse Matrices and Change of Basis}

\item Find the inverse of each matrix, or state that it is singular.
\begin{enumerate}
  \item $\begin{pmatrix} 3 & 1 \\ 2 & 1 \end{pmatrix}$

  \item $\begin{pmatrix} 2 & 4 \\ 1 & 2 \end{pmatrix}$

  \item $\begin{pmatrix} 1 & 0 & 2 \\ 0 & 1 & 0 \\ 0 & 0 & 1 \end{pmatrix}$
\end{enumerate}

\item For the linear B\'{e}zier case ($n=1$), the conversion matrix between the
Bernstein basis $\{1-u, u\}$ and the monomial basis $\{1, u\}$ is
\[
  M = \begin{pmatrix} 1 & 0 \\ -1 & 1 \end{pmatrix}.
\]
\begin{enumerate}
  \item Verify that $[B_{0,1}(u),\, B_{1,1}(u)] = [1, u]\, M$, i.e., that
        $(1-u) = 1 \cdot 1 + u \cdot (-1)$ and $u = 1 \cdot 0 + u \cdot 1$.
  \item Find $M^{-1}$.
  \item The power basis curve $C(u) = \mathbf{a}_0 + \mathbf{a}_1 u$ has $\mathbf{a}_0 = (1,0)$,
        $\mathbf{a}_1 = (2,1)$.
        Using $[P_i] = M^{-1}[\mathbf{a}_i]$, find the equivalent B\'{e}zier control points.
\end{enumerate}

% ---------------------------------------------------------
\subsubsection*{Independence and Basis}

\item Determine whether each set of vectors is linearly independent.
If dependent, express one vector as a combination of the others.
\begin{enumerate}
  \item $\{(1, 2), (3, 6)\}$
  \item $\{(1,0,1), (0,1,1), (1,1,0)\}$
  \item $\{(1,1,0), (0,1,1), (1,0,-1)\}$
\end{enumerate}

\item The Bernstein polynomials of degree 2 are
$B_{0,2}(u) = (1-u)^2$, $B_{1,2}(u) = 2u(1-u)$, $B_{2,2}(u) = u^2$.
\begin{enumerate}
  \item Write each as a linear combination of $\{1, u, u^2\}$.
  \item Verify that $B_{0,2} + B_{1,2} + B_{2,2} = 1$ (partition of unity).
  \item Write the $3 \times 3$ matrix $M$ whose rows are the monomial
        coefficients of $B_{0,2}$, $B_{1,2}$, $B_{2,2}$.
\end{enumerate}

% ---------------------------------------------------------
\subsubsection*{Cross Product and Surface Normals}

\item Compute the cross product $\mathbf{v} \times \mathbf{w}$ and verify orthogonality
by checking $\mathbf{v} \cdot (\mathbf{v} \times \mathbf{w}) = 0$.
\begin{enumerate}
  \item $\mathbf{v} = (1, 0, 0)$, $\mathbf{w} = (0, 1, 0)$
  \item $\mathbf{v} = (1, 2, 3)$, $\mathbf{w} = (-1, 0, 1)$
  \item $\mathbf{v} = (2, -1, 0)$, $\mathbf{w} = (0, 3, -1)$
\end{enumerate}

\item A parametric surface patch has $S_u = (1, 0, 2)$ and $S_v = (0, 1, -1)$
at a point.
\begin{enumerate}
  \item Compute $S_u \times S_v$.
  \item Find the unit normal vector $N = (S_u \times S_v)/\|S_u \times S_v\|$.
  \item Find the area of the parallelogram spanned by $S_u$ and $S_v$.
\end{enumerate}

% ---------------------------------------------------------
\subsubsection*{Linear and Affine Transformations}

\item For each matrix $A$, describe the geometric effect of the transformation
$T(\mathbf{x}) = A\mathbf{x}$ on $\mathbb{R}^2$, and compute $T(1, 0)$ and $T(0, 1)$.
\begin{enumerate}
  \item $A = \begin{pmatrix} 2 & 0 \\ 0 & 2 \end{pmatrix}$

  \item $A = \begin{pmatrix} 0 & -1 \\ 1 & 0 \end{pmatrix}$

  \item $A = \begin{pmatrix} 1 & 0 \\ 0 & -1 \end{pmatrix}$
\end{enumerate}

\item Find the standard matrix of the linear transformation $T : \mathbb{R}^2 \to \mathbb{R}^2$
that rotates vectors counter-clockwise by $45^\circ$.
\emph{Hint: find $T(\mathbf{e}_1)$ and $T(\mathbf{e}_2)$ using trigonometry.}

\item The affine transformation $\Phi(\mathbf{r}) = A\mathbf{r} + \mathbf{v}$ has
\[
  A = \begin{pmatrix} \cos 30^\circ & -\sin 30^\circ \\ \sin 30^\circ & \cos 30^\circ \end{pmatrix},
  \qquad
  \mathbf{v} = \begin{pmatrix} 3 \\ 1 \end{pmatrix}.
\]
The quadratic B\'{e}zier curve has control points $P_0 = (0,0)$, $P_1 = (1,1)$, $P_2 = (2,0)$.
By affine invariance, the transformed curve has control points
$\Phi(P_0)$, $\Phi(P_1)$, $\Phi(P_2)$.
Compute the three new control points.

\item A $3$-point homogeneous control point is $\tilde{P} = (w \cdot x, w \cdot y, w)$.
\begin{enumerate}
  \item For $P = (2, 3)$ with weight $w = 2$, write the homogeneous point $\tilde{P}$.
  \item Apply the perspective map $H$ to recover $P$ from $\tilde{P}$.
  \item For $\tilde{P}_0 = (1,0,1)$, $\tilde{P}_1 = (1,1,1)$, $\tilde{P}_2 = (0,2,2)$
        (the homogeneous circular arc from Piegl \& Tiller Example 1.13),
        compute the de Casteljau midpoint
        $\tilde{C}^w(1/2)$ using the formula
        $\tilde{C}^w = (1-u)^2\tilde{P}_0 + 2u(1-u)\tilde{P}_1 + u^2\tilde{P}_2$
        at $u = 1/2$, then project via $H$ to find $C(1/2)$.
\end{enumerate}

\end{enumerate}

\clearpage
\fi
